\section{Discussion}
Across five Perturb-seq screens analysed under a fixed plan, transcriptional responses to genes absent from training were predictable only to a small extent by the models we evaluated, and only in the two large screens. There, models beat the mean of other perturbations reproducibly but removed about 5\% of its squared error. Their perturbation-specific correlations of about 0.2 compare with technical reproducibility limits of 0.72--0.85. In the discovery screen, whose metrics are well calibrated and whose responses are highly reproducible, no model improved on the mean at all. The difference tracks the number of training perturbations, and performance in the large screens rose steadily with it. Among gene pairs selected for interactions, responses were well described by a scaled sum of the single responses, and learned models of the remainder were not distinguishable from that one-parameter rule.

\paragraph{Scope of the conclusions.} Our model set was deliberately simple: linear and nearest-neighbour baselines and small neural-network ensembles, using co-expression, target baseline expression or Gene2vec features. We did not evaluate graph neural networks such as GEARS \citep{roohani2024} or pretrained single-cell foundation models \citep{cui2024}. The results show how much predictive information these representations carry with these model classes, which is what limits the models tested here. They do not show that Perturb-seq data are fundamentally insufficient for stronger models, and they do not rank the field's methods. The protocol, however, applies to any model, and the reproducibility limits and calibration checks do not depend on the model at all.

\paragraph{Relation to the baseline debate.} The results bear on whether deep models beat simple baselines \citep{ahlmanneltze2025,csendes2025,vinastorne2025,miller2026}. We confirm the core of \citet{miller2026}'s argument: MSE was poorly calibrated in four of five screens, ranking half-data replicates below the mean for most perturbations. We also extend the observation of \citet{csendes2025} that common benchmarks have low perturbation-specific variance. With batch-disjoint replicates, between a quarter and two-thirds of perturbations in the large CRISPRi screens, and almost all in the neuron screens, have specific responses that do not reproduce. Remedies depend on the reference, though. Perturbation-specific correlation looks perfectly calibrated against the mean, yet a zero prediction scores well on it whenever perturbations vary mostly in magnitude, or whenever the training mean is dominated by control-sampling noise. Once noise is quantified and more than one uninformative control is used, the answer for the models tested is neither ``learned models fail'' nor ``learned models succeed''. They extract a real but small signal when training sets are large.

\paragraph{What the models learn.} The feature ablation helps interpret that signal. Much of it is available to a nonlinear model from the target's own expression level and detection rate in unperturbed cells. These are generic properties of the gene, and they relate to response magnitude only weakly and non-monotonically. Signal of this kind can help rank perturbations, but it says little about which downstream genes change. Co-expression across control cells adds a component that even a linear model can use, consistent with broad co-regulated programmes in which genes that vary together at baseline also share responses to perturbation. Two observations argue that this signal is generic rather than mechanistic: specific correlations stay far below the technical limits, and a measured response of the same gene in another cell line was only about as informative as the models. Most of a perturbation's specific transcriptional response therefore appears to depend on cellular context and is not recoverable from these representations. Generic pretrained embeddings did not close the gap: Gene2vec, built from co-expression across 984 datasets, did not beat co-expression measured in each screen's own control cells.

\paragraph{Implications for prioritising experiments.} The practical value is correspondingly limited, with three qualifications. First, in the large screens model rankings found strong responders at about 1.7 times the chance rate, a modest but usable gain. Second, for pairs whose singles are measured, nearly all of the predictable signal comes from the singles, so a scaled sum is a strong, transparent default. Third, and most important for practice, prediction intervals were reliable on average but covered only 35--55\% of the most strongly responding genes. Stratifying calibration by predicted magnitude did not fix this, because the models do not anticipate which genes will respond. Uncertainty estimates for the genes that matter most should not be relied on with these models.

\paragraph{Methodological lessons.} Several points generalise. Training and evaluation targets should not share control cells. With few controls, a shared control population couples every response to the same noise and changes which baseline looks best: the training mean was worse than predicting no change in both neuron screens. Reproducibility limits should come from disjoint batches and be labelled as technical limits. Here the batch-disjoint and split-cell limits agreed closely except in the Tian CRISPRa screen, where the batch-disjoint limit was lower (\numCeilBatchTianCrispra{} versus \numCeilHalfTianCrispra). Metric calibration should be checked against more than one uninformative prediction. Composite decision rules should report both of their components: a model can be better for most perturbations yet not on average. Finally, conclusions drawn from a single dataset can be wrong in either direction. The discovery screen suggested that no model beats the mean, and in larger screens that conclusion was refuted, if only modestly.

\paragraph{Limitations.} The study covers five screens from two laboratories. Two candidate screens were excluded for label problems in the harmonised files that study-specific curation could have resolved. Genome-wide screens and other modalities were not analysed. Batches are technical lanes within each experiment, so the reproducibility limits omit biological replicate variation, such as independent cultures or repeated experiments. Our results therefore describe predictability within an experiment, and performance against an independent repeat of a screen would be lower. The pair analysis rests on one screen whose pairs were selected for interactions. All discovery-set baselines and metrics were defined after seeing discovery data. Several analyses were defined after seeing confirmation data, including the zero-prediction calibration control, the reproducible-subset sensitivity analysis, the feature ablation and the stratified intervals; these are labelled post hoc throughout. The analysis plan was not registered with a third party.
